\section{The harmonic coefficient and its polylogarithmic representations}
\label{sec:harmonic}

The coefficient $\eta(a)$ in \eqref{ord:etadef} is not an undefined
remainder. It has an absolutely convergent Stieltjes series, belongs
to an established harmonic-zeta convention, and admits a convergent
Mellin representation with an elementary kernel when $a=1$.
This section also records exact difference and derivative identities.

\subsection{Strict and inclusive harmonic conventions}

For $n\ge0$ define the strict shifted harmonic sum
\[
 H_n^{(s)}(a)=\sum_{k=0}^{n-1}(k+a)^{-s},\qquad H_0^{(s)}(a)=0.
\]
It is entire in $s$ for fixed $n$. Its spectral derivatives are
finite logarithm-weighted harmonic sums. Summing first over the
inner index gives
\begin{equation}\label{harm:strict}
 Z_2(s,1;a)=\sum_{n\ge0}\frac{H_n^{(1)}(a)}{(n+a)^s}.
\end{equation}
The canonical decomposition with the second deviation zero yields
\begin{equation}\label{harm:etaLaurent}
 Z_2(1+u,1;a)=\frac1{u^2}+\frac{\gamma_0(a)}u+
 B_2(a)+\eta(a)u+O(u^2).
\end{equation}
Indeed the regular part is $R_2(u,0;a)$ and its first coefficient is
exactly \eqref{ord:etadef}.

Karg\i n et al. use the inclusive harmonic Hurwitz function
\cite{Kargin}
\[
 \zeta_H(s,a)=\sum_{n\ge0}(n+a)^{-s}
                           \sum_{k=0}^{n}(k+a)^{-1}.
\]
The relation and the coefficient conversion are
\begin{align}
 \zeta_H(s,a)&=Z_2(s,1;a)+\zeta(s+1,a),\label{harm:inclusive}\\
 \gamma_H(0,a)&=\frac{\gamma_0(a)^2+\zeta(2,a)}2,\notag\\
 \eta(a)&=\boxed{-\gamma_H(1,a)-\zeta'(2,a).}\label{harm:conversion}
\end{align}
Here $\gamma_H(m,a)$ is normalized by the coefficient
$(-1)^m\gamma_H(m,a)/m!$ in the regular Laurent series of
$\zeta_H(1+u,a)$. This fixes both the strict/inclusive difference and
the sign of the first coefficient.

\subsection{Shift and parameter-derivative identities}

\begin{proposition}[Transport of the ordering coefficient]
\label{harm:transport}
For every $a>0$,
\begin{equation}\label{harm:shift}
 \boxed{\eta(a+1)-\eta(a)=\frac{\gamma_1(a+1)}a.}
\end{equation}
Writing $\partial_1,\partial_2$ for the spectral derivatives of
$Z_2$, one also has the convergent identity
\begin{align}\label{harm:derivative}
 \eta'(a)={}&-Z_2(2,1;a)-\partial_1Z_2(2,1;a)
                 +\partial_2Z_2(2,1;a)\notag\\
 &+\gamma_1(a)\zeta(2,a)+\zeta'(3,a).
\end{align}
All three double sums on the right converge absolutely, including
their indicated logarithmic insertions.
\end{proposition}

\begin{proof}
Removing the terms with smallest index zero gives
\[
 Z_2(s,1;a)-Z_2(s,1;a+1)=a^{-1}\zeta(s,a+1).
\]
The coefficient of $s-1$ is
$\eta(a)-\eta(a+1)=-\gamma_1(a+1)/a$, proving
\eqref{harm:shift}.

In an absolute-convergence domain differentiation in the common
shift gives
\[
 \partial_a Z_2(s,1;a)=-sZ_2(s+1,1;a)-Z_2(s,2;a).
\]
Use the stuffle identity
\[
 Z_2(s,2;a)=\zeta(s,a)\zeta(2,a)-Z_2(2,s;a)-\zeta(s+2,a)
\]
and continue near $s=1$. Taking the coefficient of $s-1$ in
\eqref{harm:etaLaurent} gives \eqref{harm:derivative}.
Convergence follows from $H_n^{(1)}(a)=O(1+\log n)$ and the extra
outer square; spectral derivatives add only logarithmic factors.
\end{proof}

For integer shifts, repeated use of \eqref{harm:shift} yields the
finite identity
\[
 \eta(a+N)=\eta(a)+\sum_{k=0}^{N-1}
                         \frac{\gamma_1(a+k+1)}{a+k}.
\]
These transport identities are consistent with the harmonic Hurwitz
shift recurrence; they are included to make the ordering coefficient
usable without a change of convention.

\subsection{A convergent Lerch--polylogarithm Mellin formula}

For $0<z<1$, let $\Phi(z,s,a)=\sum_{n\ge0}z^n/(n+a)^s$.
The harmonic generating function is
\begin{equation}\label{harm:gen}
 \sum_{n\ge0}H_n^{(1)}(a)z^n=\frac{z\Phi(z,1,a)}{1-z}.
\end{equation}
It follows by interchanging two absolutely convergent sums. Hence
\begin{equation}\label{harm:Mellin}
 \Gamma(s)Z_2(s,1;a)=\int_0^\infty t^{s-1}K_a(t)\,dt,
 \quad K_a(t)=\frac{e^{-(a+1)t}\Phi(e^{-t},1,a)}{1-e^{-t}},
 \quad \Re s>1.
\end{equation}

Put $h_a=\gamma_0(a)-\gamma$ and define
\begin{align}\label{harm:J}
 J_a(s)={}&\int_0^1t^{s-1}
       \left(K_a(t)-\frac{-\log t+h_a}{t}\right)\,dt\notag\\
 &+\int_1^\infty t^{s-1}K_a(t)\,dt.
\end{align}
The kernel subtraction is $O(1+|\log t|)$ at zero. Indeed, put
$d_n=(n+a)^{-1}-(n+1)^{-1}$. Then $d_n=O((n+1)^{-2})$,
$\sum d_n=h_a$, and
\[
 \sum_{n\ge0}\frac{|1-e^{-nt}|}{(n+1)^2}
        =O\bigl(t(1+|\log t|)\bigr).
\]
Comparing $\Phi(e^{-t},1,a)$ with $\Phi(e^{-t},1,1)$ proves the
asserted estimate, locally uniformly for $a>0$.
Thus $J_a$ is holomorphic for
$\Re s>0$, with all fixed spectral derivatives obtained by inserting
powers of $\log t$. We have the exact continuation
\begin{equation}\label{harm:Jidentity}
 \Gamma(s)Z_2(s,1;a)=\frac1{(s-1)^2}+
                 \frac{h_a}{s-1}+J_a(s).
\end{equation}

\begin{proposition}[A normally convergent integral for $\eta$]
\label{harm:etaIntegral}
Let $g_0=\gamma_0(a)$. Then
\begin{align}\label{harm:etaintegral}
 \eta(a)={}&J_a'(1)+\gamma B_2(a)
 -\frac{\gamma^2+\zeta(2)}2g_0
 +\frac{\gamma^3}6+\frac{\gamma\zeta(2)}2+\frac{\zeta(3)}3,\\
 J_a(1)={}&\frac{(g_0-\gamma)^2+\zeta(2)-\zeta(2,a)}2.\label{harm:Jone}
\end{align}
\end{proposition}
\begin{proof}
Multiply \eqref{harm:etaLaurent} by the Taylor expansion of
$\Gamma(1+u)$ through cubic order and compare the constant and linear
coefficients with \eqref{harm:Jidentity}. The expansion of
$\log\Gamma(1+u)$ gives
\[
 [u^2]\Gamma(1+u)=\frac{\gamma^2+\zeta(2)}2,\quad
 [u^3]\Gamma(1+u)=-\frac{\gamma^3}6-
                  \frac{\gamma\zeta(2)}2-\frac{\zeta(3)}3,
\]
which proves the two formulas.
\end{proof}

At $a=1$, $\Phi(z,1,1)=\Li_1(z)/z=-\log(1-z)/z$ and
\[
 K_1(t)=-\frac{\log(1-e^{-t})}{e^t-1},\qquad J_1(1)=0.
\]
The integral formula reduces to
\begin{align}\label{harm:etaone}
 \eta(1)={}&\frac{\gamma^3}6-\frac{\gamma\zeta(2)}2+
                \frac{\zeta(3)}3\notag\\
 &+\int_0^1\log t\left(
      -\frac{\log(1-e^{-t})}{e^t-1}+\frac{\log t}{t}\right)dt
 +\int_1^\infty\frac{-\log t\log(1-e^{-t})}{e^t-1}\,dt.
\end{align}
Both integrals converge ordinarily. Numerically,
$\eta(1)=0.5367870243584185675036078843488117\ldots$.
The numerical section compares this integral with Cauchy extraction
from an independently continued ordered double sum.

All higher strict harmonic coefficients follow from the same kernel;
they do not require a separate continuation. Put
\[
 e_j=[u^j]\frac1{\Gamma(1+u)},\qquad
 C_m(a)=[u^m]R_2(u,0;a).
\]
Multiplying \eqref{harm:Jidentity} by $1/\Gamma(1+u)$ gives, for
every integer $m\ge0$,
\begin{align}\label{harm:allcoefficients}
 C_m(a)&=e_{m+2}+h_a e_{m+1}
       +\sum_{j=0}^m\frac{e_{m-j}}{j!}J_a^{(j)}(1),\\
 C_m(a)&=\frac{(-1)^m\gamma_H(m,a)-\zeta^{(m)}(2,a)}{m!}.
                  \label{harm:allconversion}
\end{align}
The second formula is \eqref{harm:inclusive} coefficient by
coefficient. Each $J_a^{(j)}(1)$ is an ordinarily convergent integral
obtained by inserting $\log^j t$ in \eqref{harm:J}. Thus the complete
strict harmonic Stieltjes hierarchy has a single convergent kernel,
with reciprocal-Gamma coefficients supplying its finite normalization.

\subsection{Polylogarithmic primitives and harmonic coefficients}

For $|z|<1$, the strict multiple polylogarithm
\[
 \Li_{s_1,\ldots,s_d}(z)=
 \sum_{n_1>\cdots>n_d\ge1}
       \frac{z^{n_1}}{n_1^{s_1}\cdots n_d^{s_d}}
\]
has the ordinary normalized Euler primitive
\begin{equation}\label{harm:primitive}
 \int_0^z\Li_{s_1,s_2,\ldots,s_d}(w)\,\frac{dw}{w}
       =\Li_{s_1+1,s_2,\ldots,s_d}(z).
\end{equation}
All fixed spectral derivatives commute with this identity on compact
subdisks. Their inner coefficients are the unshifted logarithm-weighted
harmonic sums that also occur in the remainder construction
\eqref{ord:Ecoeff}. At all indices one,
\begin{equation}\label{harm:allones}
 \Li_{\{1\}^d}(z)=\frac{[-\log(1-z)]^d}{d!},\qquad
 \sum_{d\ge1}y^{d-1}\Li_{s,\{1\}^{d-1}}(z)
 =\sum_{n\ge1}\frac{z^n}{n^s}
             \frac{\Gamma(n+y)}{\Gamma(n)\Gamma(1+y)}.
\end{equation}
The second identity follows by taking elementary symmetric
polynomials of $1,1/2,\ldots,1/(n-1)$; the Gamma quotient is the
polynomial $\prod_{j=1}^{n-1}(1+y/j)$, including at its apparently
exceptional parameters. One may first take $y$ near zero and then
continue in $y$. These are classical
polylogarithmic and harmonic generating identities, recorded here
to give a concrete derivative and antiderivative dictionary for the
Laurent coefficients. The extension at the singular endpoint is the
specified spectral calculus proved above; the endpoint is not obtained
by inserting $z=1$ into a divergent series.
